\section*{Kritische Analyse des Geschäftsmodells der Russmedia GmbH}

\subsection*{Abgrenzung und Geschäftsmodell im Business Model Canvas}

Russmedia GmbH wird in dieser Analyse als Gesamtunternehmen betrachtet, jedoch ohne die 2026 vollzogene Übernahme der Moser Holding zu berücksichtigen. Zur Russmedia GmbH gehören verschiedene Medienangebote in Print und Digital, Apps, Antenne Vorarlberg und Ländle TV sowie Anzeigen-, Druck-, Projekt-, Veranstaltungs- und Kooperationsleistungen. Das Unternehmen verbindet damit mehrere Geschäftsfelder und Zielgruppen. Das Business Model Canvas (BMC) macht die wesentlichen Bausteine sichtbar, die Zusammenfassung in einem Canvas verdeckt jedoch, dass sich Erlöslogik, Kosten und Kundenerwartungen je nach Geschäftsfeld unterscheiden.

\paragraph*{Kundensegmente.}
Auf der Publikumsseite richtet sich Russmedia an Leser:innen und Nutzer:innen, darunter zahlende Abonnent:innen und nicht zahlende Nutzer:innen. Auf der Geschäftskundenseite zählen Unternehmen zu den Kund:innen, die Werbeanzeigen, Promotion-Texte, Stellenangebote, Beilagen, Druckaufträge, Kundenprojekte oder Kooperationen beauftragen. Politische Akteure treten als Anzeigenkund:innen und teilweise als Fördergeber auf. Dabei sind Nutzer:innen nicht immer zugleich zahlende Kund:innen: Für das digitale Anzeigengeschäft ist die Reichweite wichtig, während Abonnements unmittelbare Erlöse erzielen. Besonders relevant sind daher Abonnent:innen, Wirtschaft und Politik. Eine zentrale Herausforderung liegt in der Demographie, denn zahlungsbereite Kund:innen werden älter, während ein Erreichen der unter 40-Jährigen dank kostenloser Alternativen schwierig ist.

\paragraph*{Nutzenversprechen.}
Für das Publikum besteht der Kernnutzen darin, zu erfahren, was in Vorarlberg und der Welt geschieht. Russmedia bietet aktuelle Nachrichten, einen Überblick über das Geschehen und teilweise exklusive Inhalte in Bild und Ton. Die regionale Verankerung, die Markenbekanntheit und die Nähe zur Region stützen dieses Versprechen. Wer in Vorarlberg eine breite Öffentlichkeit erreichen möchte, findet in den Angeboten des Unternehmens eine hohe regionale Reichweite. Für Werbekund:innen ist diese Reichweite daher ein wesentlicher Nutzen. Das Versprechen steht jedoch unter Druck, denn die wahrgenommene Qualität der angebotenen Medien hat laut Ansicht der Nutzer:innen nachgelassen, Preise sind gestiegen und immer mehr Inhalte werden hinter eine Paywall gestellt. Gleichzeitig konkurrieren kostenlose Angebote, insbesondere soziale Medien, um Aufmerksamkeit.

\paragraph*{Kanäle.}
Russmedia erreicht Zielgruppen über eigene Webseiten, Printprodukte, Apps, Antenne Vorarlberg und Ländle TV sowie über Instagram, Facebook, TikTok und YouTube. Veranstaltungen wie die Russ-Ausfahrt und der Russ-Preis, Sponsorings von Veranstaltungen und Sportvereinen sowie die Namenspartnerschaft von VOL.AT mit der höchsten Vorarlberger Fußballliga schaffen zusätzliche Sichtbarkeit und Kontaktpunkte. Die Kanäle werden grundsätzlich für mehrere Zielgruppen genutzt, wobei sich einzelne Veranstaltungen besonders an die Wirtschaft richten. Soziale Medien sind kostengünstig und wichtig, doch ihre große Zahl kostenloser Alternativen erschwert die Umwandlung junger Nutzer:innen in zahlende Abonnent:innen. Kostenlos angebotene Abonnements der Vorarlberger Nachrichten und der NEUE an Schulen und der Fachhochschule schaffen zwar Kontakt mit jungen Menschen, ihre langfristige Wirkung auf Nutzung und Zahlungsbereitschaft ist jedoch offen.

\paragraph*{Kundenbeziehungen.}
Das Unternehmen bietet Kundenservice, Newsletter, Apps, persönliche Betreuung und Abo-Angebote an. Eine Key-Account- und Regionalverkaufsabteilung betreut Geschäftskund:innen. Die Kundenbindung bei Abonnent:innen stützte sich bisher stark auf die Marktstellung und das reguläre Produkt. Angebote richten sich in den meisten Fällen ausschließlich an Neukund:innen, was bestehende Abonnent:innen ans Ungleichbehandung empfinden. Der Abo-Vorteilsclub soll Bestandskund:innen halten, wird aber als vom Mehrwert überschaubar wahrgenommen. Hinzu kommen jährliche Preiserhöhungen im Juli sowie Unzufriedenheit über sinkende Qualität und ausgebliebene Zustellungen. Die neu gegründete Retainment-Abteilung reagiert auf die bislang geringe Priorität systematischer Kundenbindung.

\paragraph*{Einnahmequellen.}
Als wichtigste Einnahmequellen gelten Abonnements, Anzeigen und Förderungen. Hinzu kommen unter anderem Druckaufträge, Projekte, Kooperationen, Veranstaltungen und werbliche Sonderformate. Abonnements und längerfristige Werbeverträge bieten eine gewisse Planbarkeit. Anzeigen werden je nach Ausspielkanal (Print, Online, Radio oder TV) unterschiedlich bepreist, die Preise sind teilweise verhandelbar. Digitale Werbeerlöse bleiben deutlich hinter den Abonnement- oder Print-Werbeerlösen zurück. Insgesamt scheinen die meisten Einnahmequellen rückläufig zu sein. Wachstum wird derzeit vor allem beim Radio und bei Online-Abonnements beobachtet, wobei letztere noch ein niedriges Niveau aufweisen.

\paragraph*{Schlüsselressourcen.}
Zentrale Ressourcen sind angestellte und freie Mitarbeitende, redaktionelle Kompetenz, starke und bekannte Marken, Leser:innen, Werbekund:innen sowie Kunden- und Presseagenturdaten. Hinzu kommen die technische Infrastruktur im Haus, Webseiten, Apps, Druckerei, Zustellnetz und die regionale Präsenz. Die Leserschaft ist wegen der begrenzten Größe des Marktes und der demografischen Entwicklung schwer zu ersetzen. Der zentrale Standort und die etablierten Marken bieten Vorteile. Auch das Mindset in den Redaktionen ist eine relevante Ressource, denn für neue Formate und kanalübergreifende Arbeit muss geprüft werden, wie bestehende Arbeitsweisen weiterentwickelt werden können.

\paragraph*{Schlüsselaktivitäten.}
Zu den wesentlichen Aktivitäten gehören Recherche und redaktionelle Aufbereitung, die Produktion und Zustellung der Zeitung, der Betrieb performanter Webseiten und Apps, Radioprogramm und TV-Angebote, Anzeigenverkauf, Kundenbetreuung sowie die Entwicklung neuer Werbeformate. Wichtig sind außerdem Veranstaltungen und Sponsorings, die die Marke in der Region sichtbar machen. Die redaktionellen Tätigkeiten unterscheiden sich über Kanäle hinweg grundsätzlich wenig, für die Zeitung wird jedoch ein höherer Input bei geringerem Output erwartet. Gleichzeitig werden Inhalte bereits in mehreren Kanälen genutzt. Eine bessere kanalübergreifende Planung könnte Doppelarbeit verringern, sofern die Inhalte passend für Format und Publikum aufbereitet werden.

\paragraph*{Schlüsselpartnerschaften.}
Russmedia arbeitet mit Presseagenturen, Softwareanbietern, Lieferanten für den Druck, Zusteller:innen, Backbone-Anbietern, verschiedenen Social Media Anbietern sowie mit Politik, Sportvereinen und Messeveranstaltern zusammen. Presseagenturen liefern Texte, Bilder und Grafiken. YouTube, Instagram und TikTok sind wichtige Plattformen für Video- und Social-Media-Inhalte. A1 stellt Connectivity bereit, die APA ist für die Zeitungsapps relevant. Eine besondere Abhängigkeit besteht im Printbereich vom Dienstleister PPI, der Zeitungsplanung, Produktion und Anzeigenverwaltung unterstützt. Für die digitale Infrastruktur im Haus sind A1 sowie Red Hat und NTS wichtige Partner. Veranstaltungen und Sponsorings schaffen Sichtbarkeit. Ihr direkter Beitrag zur Gewinnung neuer Abonnent:innen ist schwer zu belegen.

\paragraph*{Kostenstruktur.}
Der größte Kostenblock ist das Personal. Weitere Kosten entstehen durch Softwarelizenzen und Serviceverträge, Druckmaterial und Wartung, Presseagentur-Lizenzen, KI-Anwendungen, Investitionen in digitale Infrastruktur, Fuhrpark, Gebäude, Sponsorings und Veranstaltungen. Freie Zusteller:innen bilden einen wichtigen Teil der physischen Verteilung. Das Printgeschäft zieht zusätzliche Kosten für Druckvorstufe, Druck, Material, Vertrieb, Zustellung und Kundenservice nach sich. Bei sinkender Auflage könnten die Druckkosten zurückgehen, zugleich ist zu prüfen, wie lange die flächendeckende physische Zustellung wirtschaftlich tragbar bleibt. Einsparungen bei Papier, Recherche oder Zustellung würden allerdings direkt das Nutzenversprechen und die Kundenzufriedenheit gefährden.

\subsection*{Kritische Würdigung und Entwicklungsoptionen}

Das Geschäftsmodell verfügt über eine starke regionale Position und ein breites Portfolio an Kanälen. Print, Online, Radio, TV und Veranstaltungen ermöglichen es, unterschiedliche Nutzungssituationen abzudecken und Geschäftskund:innen Reichweite anzubieten. Die Marken und die regionale Nähe sind schwer kurzfristig aufzubauen. Dieses Potenzial wird jedoch durch eine rückläufige Entwicklung der meisten Erlösquellen und eine alternde zahlungsbereite Leserschaft begrenzt. Die hohe Reichweite allein sichert keine Abonnements, denn wenn Qualität, Zustellung und Preis-Leistungs-Verhältnis nicht überzeugen, können Nutzer:innen auf kostenlose Alternativen ausweichen. Steigende Paywall-Anteile können Erlöse je Nutzer:in erhöhen, zugleich aber Reichweite und Werbeinventar schwächen. Das Unternehmen muss daher die Balance zwischen direkter Finanzierung durch Abos und möglichst breiter Publikumsreichweite laufend prüfen.

Aus diesen Spannungen ergeben sich drei Entwicklungsoptionen:

\begin{enumerate}
  \item \textbf{Kundenbindung und Abonnements konsequent weiterentwickeln.} Die neu geschaffene Retainment-Abteilung kann Kündigungsgründe systematisch erfassen und Maßnahmen für Bestandskund:innen entwickeln. Dazu gehören nachvollziehbare Preis-Leistungs-Kommunikation, bessere Vorteile für bestehende Abonnent:innen und ein gezielter Einsatz von Test-Abos und Gewinnspielen. Exklusive regionale Geschichten mit persönlichem Zugang können das Abo-Versprechen stärken. Erfolgskriterien wären unter anderem Kündigungsquote, Verlängerungsrate und Umwandlung von Test-Abos in reguläre Abos.
  \item \textbf{Jüngere Zielgruppen mit regional relevanten Formaten erschließen.} Russmedia sollte Themen und Darstellungsformen für jüngere Menschen weiterentwickeln, ohne kurzfristigen Trends unkritisch zu folgen. Social Media kann dabei als Kontakt- und Einstiegsweg dienen. Entscheidend ist, ob daraus wiederkehrende Nutzung eigener Angebote entsteht. Denkbar sind kanalübergreifende Formate mit klar erkennbarem regionalem Mehrwert und ein schrittweiser Übergang zu digitalen Abonnements. Die kostenlose Präsenz an Schulen und der Fachhochschule könnte als Testfeld dienen, um Nutzung und spätere Zahlungsbereitschaft zu beobachten.
  \item \textbf{Produktion und Zustellung schrittweise anpassen.} Eine gemeinsame Themen- und Produktionsplanung für Print, Online, Social Media, Radio und TV kann redundante Arbeit reduzieren und Inhalte gezielter für unterschiedliche Kanäle aufbereiten. Parallel könnte Russmedia ausgewählte Gebiete oder Kundengruppen auf eine digitale Zustellung umstellen. Eine solche Veränderung sollte zunächst mit klaren Kriterien zu Kosten, Nutzung, Kundenzufriedenheit und Zustellqualität geprüft werden. Sie birgt das Risiko, insbesondere ältere Leser:innen auszuschließen, weshalb sie geeignete Übergänge und eine verlässliche digitale Nutzungserfahrung benötigt.
\end{enumerate}

Die Optionen ergänzen einander, setzen aber unterschiedliche Prioritäten. Kundenbindung und Qualität adressieren unmittelbar die Stabilität der Abos. Neue Zielgruppen schaffen eine längerfristige Perspektive für digitale Erlöse. Prozess- und Zustellungsänderungen können Kosten beeinflussen, dürfen das journalistische Produkt und die Verlässlichkeit der Leistung aber nicht schwächen. Für die Umsetzung sind daher messbare Pilotprojekte sinnvoll. Sie sollten Reichweite, Nutzung, Abonnements, Kündigungen, Kosten und Kundenzufriedenheit gemeinsam betrachten.
